\chapter{Four entrances to the subject}
\index{pluripotential theory!four frameworks}
\label{chap:four-entrances}

Pluripotential theory studies potentials whose curvature is positive and the
measures obtained by taking powers of that curvature.  Over the complex
numbers the local operator is $\ddc$ and positivity is tested on complex
lines.  Over a non-Archimedean field there are no ordinary smooth complex
coordinates, but there are several replacements: models over the valuation
ring, piecewise affine functions on skeleta, tropical differential forms, and
norms on spaces of sections.  Each replacement sees a different part of the
same geometry.

The four theories compared in this book correspond to four ways of entering
this circle of ideas.  The complex entrance supplies the images that
everything else imitates: Hermitian metrics, curvature currents, and
Monge--Amp\`ere measures.  The model entrance converts vertical algebraic
geometry over the valuation ring into potentials and atomic measures; pushed
to its limits it becomes the global variational theory of
Boucksom--Favre--Jonsson in \cref{part:bfj}.  The differential entrance
replaces holomorphic coordinates by tropical charts and builds the local
calculus of Chambert--Loir--Ducros developed in \cref{part:cld}.  The
homogeneous entrance takes the ground absolute value to be trivial, replaces
vertical divisors by filtrations and test configurations, and leads to the
Boucksom--Jonsson theory and the envelope theorem of \cref{part:bj}.  These
four entrances share a valuative substrate: a point of a non-Archimedean
analytic space records a rate of degeneration.  The complex and
transcendental constructions of \cref{part:complex} re-enter through the
first door, and \cref{part:comparisons} proves the exact comparisons.

This opening chapter introduces the four entrances and their common valuative
substrate.  It requires no earlier chapter and is meant to be read twice:
once now, for
orientation, and once after \cref{part:foundations}, when every object appearing below has a
precise definition.  Each entrance comes with one explicit computation,
chosen so that the same answer will be recomputed later inside the formal
theory; the reader who works through these small calculations will recognize
every later sign and normalization on first sight.  Three examples will
accompany us throughout the book:
a nodal degeneration, the projective toric line, and a test configuration.
They are set up completely in \cref{sec:ch01-running-examples} and revisited
in \cref{chap:model-calculations} and in the example chapters of every later
part.  The chapter closes with the master comparison diagram, stated without
proof, and with a reader's guide to the six parts.

\section{The complex entrance: metrics and curvature}
\label{sec:ch01-complex-entrance}
\index{metric!Hermitian}\index{curvature current}

Let $X$ be a compact complex manifold and let $L$ be a holomorphic line
bundle.  If $h_0$ is a smooth Hermitian metric on $L$, every other possibly
singular metric can be written, in terms of squared norms,
\[
             |s|_h^2=\rme^{-\varphi}|s|_{h_0}^2.
\]
Thus a metric is encoded by an extended-real function $\varphi$, called its
\emph{potential relative to $h_0$}.  With the normalization
\[
             \dc=\frac{\rmi}{4\pi}(\bar\partial-\partial),
             \qquad \ddc=\frac{\rmi}{2\pi}\partial\bar\partial,
\]
its curvature current is
\[
             c_1(L,h)=c_1(L,h_0)+\ddc\varphi.
\]
The metric is semipositive precisely when this current is positive.  If
$\varphi$ is locally bounded, Bedford--Taylor theory defines
$(c_1(L,h_0)+\ddc\varphi)^n$; finite-energy theory later extends this
construction to much more singular potentials.

The two theorems that shaped the subject should be kept in mind from the
start.  Bedford and Taylor showed that the nonlinear operator
$\varphi\mapsto(\ddc\varphi)^n$ makes sense for locally bounded psh functions
and is continuous along decreasing limits \cite{BT76,BT82}.  Yau's solution
of the Calabi conjecture showed that on a compact K\"ahler manifold the
Monge--Amp\`ere equation prescribing this measure is solvable \cite{Y78}.
The modern global theory combines the two: the variational method of
\cite{BBGZ13} solves degenerate Monge--Amp\`ere equations by maximizing an
energy functional over a compact set of singular potentials, and the
finite-energy formalism of \cite{BEGZ10} explains on which singular
potentials the measure still has full mass.  Every one of these steps has a
non-Archimedean counterpart in this book: compactness of psh functions,
envelopes, energy, and a variational existence theorem.  A systematic
complex reference is \cite{GZ17}; the exact complex statements used later,
with our normalizations, are collected in \cref{chap:complex-prerequisites}
and \cref{app:complex-reference}.

Three features should be retained.  First, a reference metric turns metrics
into functions, but the metric itself is intrinsic.  Second, positivity is a
curvature condition.  Third, the Monge--Amp\`ere operator is nonlinear but is
continuous along suitable monotone approximations.  All three statements
have non-Archimedean analogues.

\begin{remark}
\label{rem:ch01-curvature-normalization}
Some sources use $\ddc=\rmi\partial\bar\partial/\pi$.  Every comparison with
complex geometry in this book uses the normalization displayed above.  The
operators $\dprimeddprime$ of Chambert--Loir--Ducros have their own natural
normalization; factors are recorded when the two calculi meet.
\end{remark}

\subsection{A model-case computation on the Riemann sphere}
\label{sec:ch01-equilibrium}
\index{equilibrium measure}\index{metric!maximum}

The following computation is the complex model case for the entire book.  It
produces, by an elementary jump computation, the curvature measure of the
most degenerate natural metric on $\cO(1)$, and its answer is the exact
complex shadow of the non-Archimedean measures computed in
\cref{chap:model-calculations}.

\begin{lemma}\label{lem:ch01-equilibrium}
On $\mathbf C$, let
\[
   u(z)=\max\{0,\log|z|^2\}=\log\max\{1,|z|^2\}.
\]
Then $\ddc u$ is the Haar probability measure on the unit circle,
\begin{equation}\label{eq:ch01-equilibrium}
   \ddc u=\frac{1}{2\pi}\,\rd\theta
   \qquad\text{on }\{|z|=1\}.
\end{equation}
\end{lemma}
\begin{proof}
The function $u$ is continuous, vanishes on the closed unit disc, and equals
$2\log r$ in polar coordinates $z=r\rme^{\rmi\theta}$ outside it, so
$\ddc u$ is well defined as a current.  Fix a test function
$g\in C^\infty_c(\mathbf C)$ and choose $R>1$ with
$\operatorname{supp}g\subset\{|z|<R\}$.  Writing $\lambda$ for Lebesgue
measure on $\mathbf C\simeq\mathbf R^2$, the normalization
$\ddc=(\rmi/2\pi)\partial\bar\partial$ gives
$\ddc g=(4\pi)^{-1}(\Delta g)\,\lambda$, whence
\[
   \langle \ddc u,g\rangle
   =\int_{\mathbf C}u\,\ddc g
   =\frac{1}{4\pi}\int_{\mathbf C}u\,\Delta g\,\rd\lambda .
\]
The integrand vanishes on the unit disc because $u$ does.  On the annulus
$A=\{1\leq|z|\leq R\}$ the function $u=2\log r$ is harmonic, and Green's
second identity gives
\[
 \int_A\bigl(u\,\Delta g-g\,\Delta u\bigr)\rd\lambda
 =\int_{\partial A}\Bigl(u\,\frac{\partial g}{\partial n}
     -g\,\frac{\partial u}{\partial n}\Bigr)\rd s,
\]
with outward normal $n$ and arclength $s$.  Here $\Delta u=0$ on $A$, and on
the outer circle $g$ and its derivatives vanish.  On the inner circle
$\{|z|=1\}$ the outward normal of $A$ points toward the origin, so
$\partial/\partial n=-\partial/\partial r$ there; moreover $u=0$ on this
circle, which kills the first boundary term, while
$\partial u/\partial r=2$ at $r=1$.  Therefore
\[
 \int_{\mathbf C}u\,\Delta g\,\rd\lambda
 =\int_{\{|z|=1\}}2g\,\rd s ,
\]
and therefore
\[
 \langle \ddc u,g\rangle
 =\frac{2}{4\pi}\int_0^{2\pi}g(\rme^{\rmi\theta})\,\rd\theta
 =\frac{1}{2\pi}\int_0^{2\pi}g(\rme^{\rmi\theta})\,\rd\theta .
\]
This is \eqref{eq:ch01-equilibrium}.  The measure is positive of total mass
$1$.
\end{proof}

The mechanism deserves as much attention as the answer: away from the corner
locus the potential is harmonic and contributes nothing, and the entire
curvature is the jump of the normal derivative across the interface.  This
is precisely the mechanism of the graph Laplacian
\eqref{eq:graph-laplacian}, where the mass at a point of a metric graph is
the sum of outgoing slopes.  The complex and non-Archimedean theories
compute corners in the same way.

\begin{example}[The maximum metric on the hyperplane bundle]
\label{ex:ch01-equilibrium-metric}
On $X=\mathbf P^1_{\mathbf C}$ with homogeneous coordinates $[Z_0:Z_1]$,
define a continuous metric on $\cO(1)$ by
\[
   \|Z_0\|([Z_0:Z_1])=\frac{|Z_0|}{\max\{|Z_0|,|Z_1|\}}.
\]
On the chart $z=Z_1/Z_0$ the squared-norm weight is
\[
   -\log\|Z_0\|^2=\log\max\{1,|z|^2\}=u(z),
\]
so \cref{lem:ch01-equilibrium} gives
\[
   c_1\bigl(\cO(1),\|\cdot\|\bigr)
   =\frac{1}{2\pi}\,\rd\theta
   \qquad\text{on the circle }|Z_1|=|Z_0|,
\]
and the chart at infinity, where the weight is
$\max\{0,\log|w|^2\}$ with $w=Z_0/Z_1$, confirms that no mass hides at
$z=\infty$.  The total mass is $1=\deg\cO(1)$, and the measure is the
equilibrium measure of the closed unit disc.  Two remarks calibrate this
example for later use.

First, the metric is a uniform limit of smooth metrics: the weights
\[
   u_m(z)=\frac1m\log\bigl(1+|z|^{2m}\bigr)
\]
are smooth, satisfy $0\leq u_m-u\leq m^{-1}\log2$, and their curvature forms
converge weakly to \eqref{eq:ch01-equilibrium}.  Concentration of smooth
positive curvature onto a corner locus is not a pathology; it is how model
geometry appears as a limit, and the toric version of the same computation
occurs in \cref{sec:toric-convex}.

Second, over a non-Archimedean field the very same formula defines the
standard model metric of
\cref{sec:berkovich-p1}, whose curvature \eqref{eq:p1-standard-measure} is
the Dirac mass $\delta_{\zeta_{0,1}}$ at the Gauss point.  The unit circle,
an orbit of the circle action, collapses to a single nonclassical point.
The hybrid spaces of \cref{chap:hybrid-spaces} make this collapse a genuine
limit of spaces and measures.  Finally, note the factor of two: complex
potentials are weights of squared norms, non-Archimedean potentials are
weights of norms, and the conversion is displayed whenever the two meet
(\cref{app:normalizations}).
\end{example}

\section{The valuative substrate: points as rates of degeneration}
\label{sec:ch01-valuative-entrance}
\index{semivaluation}\index{degeneration!rate of}

Let $(K,|\cdot|)$ be a complete non-Archimedean field.  Our additive sign
convention is
\[
                        v(a)=-\log|a|.
\]
Consequently $v(ab)=v(a)+v(b)$ and
$v(a+b)\geq\min\{v(a),v(b)\}$.  If $A$ is a finitely generated $K$-algebra,
a point $x$ of the analytification $(\Spec A)^{\An}$ is a multiplicative seminorm
$|\cdot|_x$ on $A$ extending $|\cdot|$ on $K$.  Equivalently it is a real
semivaluation
\[
       v_x:A\longrightarrow\mathbf R\cup\{+\infty\},
       \qquad v_x(f)=-\log|f(x)|.
\]
Allowing $+\infty$ is essential: the set
$\{f\in A\mid v_x(f)=+\infty\}$ is the support of the point and need not be
zero.

The slogan ``a point is a rate of degeneration'' becomes literal on a model.
Assume for the moment that $K$ is discretely valued, choose a uniformizer
$\pi$ with $v(\pi)=1$, and let $\mathcal X$ be a normal model of a proper
$K$-variety $X$.  Write
\[
                    \mathcal X_s=\sum_E b_EE.
\]
Every irreducible component $E$ determines a valuation
\[
                  v_E(f)=\frac{1}{b_E}\operatorname{ord}_E(f)
                  \quad(f\in K(X)^\times),
\]
whose restriction to $K$ is $v$.  The corresponding analytic point records
the vanishing rate of every rational function along $E$.  Blowing up the
special fibre produces more such points; their closure is far richer than
the set of closed algebraic points.

\begin{example}[The Gauss points of the affine line]
\label{ex:ch01-gauss-points}
For $r\in\mathbf R$, define on $K[T,T^{-1}]$
\[
 \left|\sum_i a_iT^i\right|_r=\max_i |a_i|\rme^{-ir},
 \qquad
 v_r\left(\sum_i a_iT^i\right)=\min_i\{v(a_i)+ir\}.
\]
The ultrametric inequality is immediate.  For multiplicativity, filter the
Laurent polynomial ring by the displayed weights.  Its associated graded
ring is
$(\operatorname{gr}_vK)[\tau,\tau^{-1}]$, with
$\deg\tau=r$, and is a domain.  The initial form of a product is therefore
the product of the initial forms, so the least weights add.  This argument
works for every real $r$; no rescaling into the value group is required.
The parameter
$r=v_r(T)$ is a tropical coordinate.  Laurent polynomials therefore restrict
to piecewise affine functions of $r$.
\end{example}

The one-line formula in \cref{ex:ch01-gauss-points} already contains the
seed of everything tropical in this book.  For a fixed Laurent polynomial
$f$, the function $r\mapsto v_r(f)$ is a minimum of finitely many affine
functions, hence concave and piecewise affine with integral slopes; its
bends occur exactly where two monomials of $f$ acquire the same size, and
the collection of bends is the tropical hypersurface of $f$.  The formal
theory of such points, their completed residue fields, and their graded
reductions occupies \cref{chap:valuations}.  Their analytic topology is
built in \cref{chap:berkovich}, which distinguishes affine analytifications
from compact affinoid domains and proves compactness for proper varieties.

\subsection{A valuation on a curve: the node and its branches}
\label{sec:ch01-node-branches}
\index{nodal curve}\index{valuation!on a curve}

Valuations also probe singularities, and the simplest instructive case is a
curve singularity over a trivially valued field.  The following computation
introduces, in miniature, the dichotomy on which the sharpest theorem of
this book turns.

\begin{example}[Two branches through a node]
\label{ex:ch01-nodal-branches}
Let $k$ be an algebraically closed field with $\operatorname{char}k\neq2$,
equipped with the trivial absolute value, and let
\[
   C=\{y^2=x^2(1+x)\}\subset\mathbf A^2_k
\]
be the nodal cubic, with node at the origin $O=(0,0)$.  In $k[[x]]$ the
element $1+x$ has a square root
$\sigma(x)=1+\tfrac{x}{2}-\tfrac{x^2}{8}+\cdots$, and the two formal
branches of $C$ through $O$ are
\[
   y=x\sigma(x),\qquad y=-x\sigma(x).
\]
The normalization $\nu:\widetilde C\to C$ has two points $p_+,p_-$ above
$O$, one for each branch, and each defines a valuation on the function field
$F=k(C)$,
\[
   v_\pm=\operatorname{ord}_{p_\pm} .
\]
Both are centered at $O$.  We compute them on generators.  Along either
branch, $x$ is a local parameter, so
\[
   v_\pm(x)=1,\qquad v_\pm(y)=v_\pm\bigl(\pm x\sigma(x)\bigr)=1 ,
\]
since $\sigma(0)=1\ne0$.  Thus $v_+$ and $v_-$ agree on $x$ and on $y$, and
indeed on the whole subring $k[x]$, where both restrict to
$\operatorname{ord}_0$.  They differ on functions that see the two branches
asymmetrically:
\begin{gather*}
 v_+(y-x)=\operatorname{ord}_x\bigl(x(\sigma(x)-1)\bigr)
         =\operatorname{ord}_x\Bigl(\tfrac{x^2}{2}-\cdots\Bigr)=2,\\
 v_-(y-x)=\operatorname{ord}_x\bigl(x(-\sigma(x)-1)\bigr)=1,
\end{gather*}
because $-\sigma(0)-1=-2\neq0$.  Symmetrically $v_+(y+x)=1$ and
$v_-(y+x)=2$.  So the node carries two distinct rank-one valuations with the
same center: each branch assigns the larger value to the equation of its
own tangent line, along which it has second-order contact.

Contrast the cusp $C'=\{y^2=x^3\}$.  Its normalization
$t\mapsto(t^2,t^3)$ has a single point above the singularity, defining
one valuation $v$ with
\[
   v(x)=2,\qquad v(y)=3 ,
\]
and there is no second branch to separate: $C'$ is unibranch at the cusp.
\end{example}

Over the node, the analytification contains two distinct rays of
semivaluations $t\,v_\pm$ pinched over one closed point; over the cusp there
is only one.  This local difference is not cosmetic.  The envelope theorem
of \cref{chap:bj-envelope-statement}, the strongest foundational result of
\cref{part:bj}, holds for a projective variety exactly when the variety is
unibranch, and the node/cusp pair of \cref{ex:ch01-nodal-branches} is the
germ of both the counterexample and the positive theorem.  More generally,
valuations serve throughout the book as probes of singularities of psh
functions, in the tradition of \cite{BFJ08,JM12}.  The classical valuation
theory needed for all of this---rank, rational rank, residue transcendence
degree, and the Abhyankar inequality
(\cref{thm:abhyankar-inequality})---comes from \cite{ZariskiSamuelII} and is
developed in \cref{chap:valuations}.

\section{The model entrance: vertical geometry becomes a potential}
\label{sec:ch01-model-entrance}
\index{model function}\index{divisor!vertical}

A model $\mathcal X$ spreads $X$ over the valuation ring $K^\circ$.  A
vertical Cartier divisor $D$ on $\mathcal X$ is invisible on the generic
fibre, but it changes the model metric on the trivial line bundle. Locally,
if $D$ has equation $f$, then $f^{-1}$ is a frame of
$\mathcal O_{\mathcal X}(D)$ and its model function is
\[
                         \varphi_D(x)=-\log|f(x)|=v_x(f).
\]
Changing $f$ by a unit does not change the function on the generic analytic
space.  Tensor products of model line bundles become sums of model
functions, pullback to a higher model leaves the analytic function unchanged,
and nefness on the special fibre becomes semipositivity.

The following example carries out the entire cycle---divisor, potential,
curvature---in the smallest possible case.  It is the non-Archimedean
counterpart of \cref{lem:ch01-equilibrium} and fixes, once and for all, the
sign conventions of the notation ledger: effective vertical divisors give
nonnegative model functions, and convex corners of potentials carry positive
mass.

\begin{example}[A vertical divisor on a blown-up projective line]
\label{ex:ch01-blowup-model-function}
Let $K$ be discretely valued, $R=K^\circ$, and $v(\pi)=1$.  Start from the
model $\cX_0=\mathbf P^1_R$ of $X=\mathbf P^1_K$, with affine coordinate
$T$, and blow up the closed point $\{T=\pi=0\}$ of the special fibre to
obtain $\cX_1\to\cX_0$.  Two charts describe $\cX_1$ over
$\operatorname{Spec}R[T]$:
\[
 \text{(i) } T=\pi s\ \text{ with coordinate }s ,
 \qquad
 \text{(ii) } \pi=Tu\ \text{ with coordinates }T,u .
\]
In chart (ii) the special fibre is $\operatorname{div}(\pi)
=\{T=0\}+\{u=0\}$: the first component is the exceptional curve $E$, the
second is the strict transform $\widetilde C$ of the original special
fibre, and both occur with multiplicity one,
\[
   \cX_{1,s}=\widetilde C+E .
\]
The divisorial point $x_E$ satisfies $v_{x_E}(T)=\operatorname{ord}_E(T)=1$;
it is the Gauss point of the disc $\{v(T)\geq1\}$, while $x_{\widetilde C}$
is the Gauss point with $v(T)=0$.  The segment
\[
  \Gamma=\{\,v(T)=t,\ 0\leq t\leq1\,\}
\]
joining them is the skeleton of the model $\cX_1$.

Now let $D=E$ as a vertical Cartier divisor.  Its local equation is $T$ in
chart (ii), $\pi$ in chart (i), and a unit at every point of
$\widetilde C$ not on $E$.  Evaluating $\varphi_E(x)=v_x(f)$, where $f$ is
the local equation at the reduction of $x$, in the three regions gives the
closed formula
\begin{equation}\label{eq:ch01-blowup-model-function}
   \varphi_E=\min\bigl\{\max\{v(T),0\},\,1\bigr\}:
\end{equation}
zero where $v(T)\leq0$, equal to $v(T)$ for $0\leq v(T)\leq1$, and constant
$1$ where $v(T)\geq1$.  On the skeleton, $\varphi_E(t)=t$; off the skeleton
the function is constant along the fibres of the retraction constructed in
\cref{chap:skeleta}.  The function is nonnegative, as the effective divisor
$E$ requires.

Its curvature can be computed twice.  The Laplacian of
\eqref{eq:graph-laplacian}, the sum of outgoing slopes, gives
\[
 \Delta\varphi_E=\delta_{x_{\widetilde C}}-\delta_{x_E},
\]
with the positive mass at the convex corner $t=0$ and the negative mass at
the concave corner $t=1$.  Intersection theory on $\cX_1$ gives the same
answer: $\widetilde C$ and $E$ meet transversely in one closed point, and
$\operatorname{div}(\pi)=\widetilde C+E$ has degree zero on the proper
curve $E$, so
\[
  (E\cdot\widetilde C)=1,
  \qquad
  (E\cdot E)=-1 ,
\]
and the atomic measure of \cref{sec:metrics-cl-measure}, whose coefficient
at $x_Y$ is $b_Y\deg\bigl(\cO(E)|_Y\bigr)$, equals
\[
 1\cdot(E\cdot\widetilde C)\,\delta_{x_{\widetilde C}}
 +1\cdot(E\cdot E)\,\delta_{x_E}
 =\delta_{x_{\widetilde C}}-\delta_{x_E}.
\]
The two computations agree, as they must.
\end{example}

The example also separates two positivity notions that beginners often
merge.  The divisor $E$ is effective, and its model function is nonnegative;
but $\cO(E)$ has degree $-1$ on $E$, so the metric is \emph{not}
semipositive.  Effectivity of the divisor is a sign condition on the
potential, whereas semipositivity is nefness on the special fibre, that is,
positivity of the curvature; the two are unrelated.  Semipositive model
metrics are studied systematically in \cref{sec:metrics-semipositive}.

This is the global algebraic entrance to non-Archimedean pluripotential
theory.  Intersection theory on $\mathcal X_s$ produces a measure.  If
$\mathcal L_1,\ldots,\mathcal L_n$ are relatively nef model line bundles on
a normal model, the associated mixed Chambert--Loir measure has the form
\[
 \sum_E b_E
   \bigl(c_1(\mathcal L_1|_E)\cdots c_1(\mathcal L_n|_E)\bigr)
   \delta_{x_E},
\]
up to the normalization of the base valuation.  In
\cref{chap:metrics,chap:bfj-ma} we
derive the exact formula and prove its invariance under change of model.
The construction goes back to Zhang's semipositive metrics on arithmetic
varieties \cite{Zhang1995} and to Chambert-Loir's measures \cite{CL06}; a
modern treatment of the continuity properties is \cite{GM19}.  By
\cref{thm:cl-continuous-extension}, the Chambert--Loir measure extends
continuously for uniform convergence from semipositive model metrics to their
uniform limits.  Singular limits of such metrics are the subject of the BFJ
theory in \cref{part:bfj}.

\section{The differential entrance: tropical charts and currents}
\label{sec:ch01-differential-entrance}
\index{tropical chart}\index{superform}

Chambert--Loir and Ducros replace complex charts by tropical charts.  On an
open subset of $\mathbf R^n$, Lagerberg's calculus has two copies of the
ordinary cotangent directions and operators $\dprime$ and $\ddprime$.  A moment map
from an analytic domain to a torus yields a tropicalization map
\[
                     x\longmapsto
       \bigl(-\log|f_1(x)|,\ldots,-\log|f_n(x)|\bigr),
\]
and superforms on the tropical image pull back to real differential forms on
the analytic domain.  One can then define positive forms, currents,
integration, $\dprimeddprime$, and local Bedford--Taylor products.

The formal device is easy to state.  A superform of bidegree $(p,q)$ on an
open set $U\subset\mathbf R^n$ is a smooth form with two sets of formal
differentials,
\[
   \alpha=\sum_{|I|=p,\;|J|=q}
     \alpha_{IJ}\,\dprime u_I\wedge\ddprime u_J ,
\]
where both $\dprime u_i$ and $\ddprime u_i$ are labelled copies of the same
real coordinate differential.  The operators $\dprime$ and $\ddprime$
differentiate into the first and second slot, and for a smooth function $F$,
\[
 \dprimeddprime F
 =\sum_{i,j}\frac{\partial^2F}{\partial u_i\partial u_j}
   \,\dprime u_i\wedge\ddprime u_j ,
\]
so that positivity of $\dprimeddprime F$ in the superform sense is exactly
positive semidefiniteness of the Hessian: convex functions are the local psh
functions of the tropical world.  This calculus, its positivity cones, and
its integration theory are constructed from scratch in
\cref{chap:superforms,chap:tropical-integration}, following the revised
Chambert--Loir--Ducros manuscript \cite{CLD25} (the original is
\cite{CLD12}); superforms themselves were introduced by Lagerberg, for whom
see the notes to this chapter.

As with the complex and model entrances, one computation shows the whole
mechanism.

\begin{example}[The corner of a piecewise affine potential]
\label{ex:ch01-lagerberg-corner}
On $\mathbf R$ with coordinate $u$, consider the convex function
\[
   F(u)=\max\{0,-u\},
\]
the tropical version of the weight in \cref{lem:ch01-equilibrium}.  It is
not smooth, but it defines a current, with
$\dprimeddprime[F]$ acting on a test function
$g\in C^\infty_c(\mathbf R)$ by
\[
 \langle\dprimeddprime[F],g\rangle
 =\int_{\mathbf R}F\,\dprimeddprime g
 =\int_{\mathbf R}F(u)\,g''(u)\,\rd u ,
\]
using $\dprimeddprime g=g''\,\dprime u\wedge\ddprime u$ and the tropical
integration convention
$\int f\,\dprime u\wedge\ddprime u=\int f\,\rd u$ of
\cref{chap:tropical-integration}.  Integrating by parts twice on
$(-\infty,0]$, where $F(u)=-u$,
\[
 \int_{-\infty}^0(-u)g''(u)\,\rd u
 =\bigl[-u\,g'(u)\bigr]_{-\infty}^0
   +\int_{-\infty}^0 g'(u)\,\rd u
 =g(0).
\]
Hence
\[
   \dprimeddprime[F]=\delta_0 ,
\]
a positive measure sitting at the corner.  The computation is the exact
analogue of \cref{lem:ch01-equilibrium}: harmonicity away from the corner,
and a slope jump of total size one at it.  The same convention makes the
Laplacian \eqref{eq:graph-laplacian} of a piecewise affine function on a
skeleton a special case of $\dprimeddprime$ once skeleta are related to
tropical charts in \cref{chap:analytic-tropicalization}.
\end{example}

On an analytic space, tropical charts are provided by invertible functions:
on the annulus $\{0\leq v(x)\leq1\}$ of the running nodal degeneration
(\cref{sec:ch01-running-nodal}), the single function $x$ tropicalizes the
annulus onto the interval $[0,1]$, and superforms on the interval pull back
to smooth forms on the annulus.  Currents, the Poincar\'e--Lelong formula,
and local wedge products of \emph{continuous} psh functions are then developed
in \cref{chap:cld-currents,chap:cld-poincare-lelong,chap:cld-bedford-taylor}.
The nonlinear domain of the CLD calculus is the continuous psh class; the
complex locally bounded Bedford--Taylor theory is deliberately \emph{not}
transported, for want of the smoothing that proves it
(\cref{chap:cld-bedford-taylor}).
Two later refinements deserve early mention: the delta-forms of
Gubler--K\"unnemann \cite{GK17,GK19} and the harmonic tropicalizations of
\cite{GJR26} enlarge the sheaf of forms for cohomological purposes
(\cref{chap:delta-forms}), and on curves the CLD notion of subharmonicity is
compared with the classical potential-theoretic one in
\cref{chap:cld-examples}, following \cite{Wanner19,GR25PL,GR25Classical}.

This local calculus and the model calculus overlap, but neither should be
defined by the other.  Model metrics may have piecewise affine, hence
nonsmooth, curvature.  Conversely, a smooth form on a tropical chart need
not arise from a single algebraic model.  The comparison theorem on their
common model core is one of the main goals of the book.

\section{The homogeneous entrance: trivial absolute values}
\label{sec:ch01-homogeneous-entrance}
\index{trivial absolute value}\index{cone of valuations}

Suppose now that $k$ carries the trivial absolute value and $X$ is an
integral variety.  A point of
$X^{\mathrm{an}}$ is still a semivaluation, but there is no distinguished
number $v(\pi)=1$.  Positive rescaling acts on points:
\[
                         (t\cdot v)(f)=t\,v(f),\qquad t>0.
\]
The trivial valuation is the vertex of this cone.  Homogeneous functions,
filtrations of section rings, and test configurations replace the vertical
geometry over a fixed valuation ring.

For a smooth projective curve the cone is easy to draw.  By
\cref{prop:trivial-p1-tree}, the analytification of $\mathbf P^1_k$ is an
infinite cone of rays: for each closed point $a$ there is the ray of
semivaluations $v_{a,t}=t\operatorname{ord}_a$, $0\leq t\leq+\infty$, all
rays glued at the trivial valuation $v_{\triv}$ at $t=0$, and each ray
ending at a semivaluation whose support is the maximal ideal of $a$.  The
scaling action moves points along their rays and fixes the vertex and the
endpoints.  Functions compatible with this action are the first objects of
the trivially valued theory.

\begin{example}[A homogeneous potential on the cone]
\label{ex:ch01-cone-ideal}
Let $X=\mathbf P^1_k$ and let $\mathfrak a\subset\cO_X$ be the ideal sheaf
of the closed point $0=[1:0]$.  Every semivaluation evaluates on ideals by
$v(\mathfrak a)=\min\{v(f)\mid f\in\mathfrak a_{c(v)}\}$, where $c(v)$ is
the center; set
\[
   \varphi_{\mathfrak a}(v)=-v(\mathfrak a).
\]
With $w$ a local coordinate at $0$, one computes on every ray:
\[
 \varphi_{\mathfrak a}(v_{0,t})=-t\,\operatorname{ord}_0(w)=-t,
 \qquad
 \varphi_{\mathfrak a}(v_{a,t})=0\ \ (a\neq0),
\]
since $\mathfrak a$ is trivial away from $0$.  At the endpoint of the ray
through $0$ the value is $-\infty$: the support of $v_{0,\infty}$ contains
$w$.  The function is homogeneous,
\[
 \varphi_{\mathfrak a}(t\cdot v)=t\,\varphi_{\mathfrak a}(v),
\]
piecewise affine on every ray, and continuous on the cone with values in
$[-\infty,0]$.  This is the smallest instance of a general mechanism:
coherent ideals, and more generally b-divisors over all birational models,
produce homogeneous potentials (\cref{chap:bj-homogeneity-bdivisors}); the
two-variable version $\varphi_{(x^2,y^3)}$ is computed in
\cref{sec:trivial-valuation-cone}.

Homogeneous functions are typically unbounded.  The bounded objects of the
theory arise by truncation: the function
\[
 \max\{\varphi_{\mathfrak a},-1\},
\]
equal to $\max\{-t,-1\}$ on the ray through $0$ and to $0$ elsewhere, is
bounded and continuous, and \cref{sec:ch01-running-tc} identifies it, up to
an additive constant, with the Fubini--Study function of the running test
configuration.  Truncating at $-c$ and letting $c\to\infty$ moves curvature
mass out along the ray.  The precise energy of the limiting function is
computed only after the energy functional has been constructed; see
\cref{chap:bj-finite-energy}.  Thus no finite-energy conclusion is being
used at this introductory stage.
The interplay between homogeneous functions and their bounded truncations
is formalized by the Legendre duality of
\cref{chap:bj-homogeneity-bdivisors} and, on the complex side, by the test
curves of \cref{chap:rays-test-curves}.
\end{example}

Let $L$ be ample and $R(X,L)=\bigoplus_{m\geq0}H^0(X,mL)$.  A linearly
bounded multiplicative filtration $\mathcal F$ assigns subspaces
$\mathcal F^\lambda H^0(X,mL)$
compatible with multiplication.  Its jumping numbers behave like weights,
while the associated Fubini--Study functions are maxima of finitely many
affine valuative functions.  Test configurations give finitely generated
filtrations; general finite-energy metrics arise after completion.  This is
the Boucksom--Jonsson entrance developed in \cref{part:bj}, whose
foundational source is \cite{BJ22}, with the K-stability motivation coming
from \cite{BHJ17,BJ20}.  Its central foundational statement, the envelope
theorem, is proved in this book for arbitrary integral projective unibranch
varieties in every characteristic
(\cref{chap:bj-envelope-statement}); \cref{ex:ch01-nodal-branches} shows the
smallest obstruction when unibranchness fails.

\section{Three running examples}
\label{sec:ch01-running-examples}
\index{running examples}

Three examples recur throughout the book, one for each axis of the theory: a
degenerating family over a discrete valuation ring, a toric variety over any
non-Archimedean field, and a test configuration over a trivially valued
field.  We set them up here completely and compute, in each, one potential
and one measure.  \Cref{chap:model-calculations} redoes these computations
inside the formal theory, and the example chapters of
\cref{part:cld,part:bfj,part:bj} return to them at each new level of
generality.  A reader who keeps these three pictures in mind can translate
almost every abstract statement of the book into a statement about an
interval, a convex function, or a weight decomposition.

\subsection{A nodal degeneration}
\label{sec:ch01-running-nodal}
\index{nodal degeneration}\index{skeleton!of an annulus}

Let $R=K^\circ$ be a discrete valuation ring and consider the formal local
model
\[
                   \mathfrak U=\operatorname{Spf}
                   R\{x,y\}/(xy-\pi).
\]
Its generic fibre contains the closed annulus
$|\pi|\leq |x|\leq1$, with $y=\pi/x$.  For $r\in[0,1]$ the Gauss valuation
satisfying
\[
                   v_r(x)=r,\qquad v_r(y)=1-r
\]
defines a point of the annulus.  These points form its skeleton, an interval.
The two endpoints correspond to the two components of the nodal special
fibre, and the coordinate $r$ is integral affine.  Vertical divisors restrict
to piecewise affine functions on this interval; their Laplacians are atomic
measures at slope breaks and endpoints.

The global form of this example is the projective conic bundle
\[
             \mathcal C=\{XY=\pi Z^2\}\subset\mathbf P^2_R ,
\]
studied in \cref{sec:running-semistable}: its generic fibre is a smooth
conic, isomorphic to $\mathbf P^1_K$, its special fibre is the union of the
two lines $E_0=\{Y=0\}$ and $E_1=\{X=0\}$ crossing at the node $[0:0:1]$,
and the formal completion along the node is $\mathfrak U$ in the affine
coordinates $x=X/Z$, $y=Y/Z$.  With $s=v(x)$ as coordinate, the skeleton is
the segment $[0,1]$ from the divisorial point $v_0$ of $E_0$ to the
divisorial point $v_1$ of $E_1$.  The polarization
$\cL=\cO_{\mathbf P^2_R}(1)|_{\mathcal C}$ has degree one on each component,
and its Chambert--Loir measure is $\delta_{v_0}+\delta_{v_1}$
(\cref{prop:running-conic-measure}).

The next computation shows what happens when the model is modified, and
introduces multiplicities.

\begin{example}[Blowing up the node]
\label{ex:ch01-node-blowup}
The total space of $\mathcal C$ is regular at the node: in
$A=R[x,y]/(xy-\pi)$ the maximal ideal of the node is $(x,y)$, and
$\pi=xy\in(x,y)^2$.  Blow up $(x,y)$.  Two charts describe the result:
\[
 \text{(i) } y=xw:\quad R[x,w]/(x^2w-\pi),
 \qquad
 \text{(ii) } x=yu:\quad R[y,u]/(y^2u-\pi).
\]
In chart (i), $\operatorname{div}(\pi)=\operatorname{div}(x^2w)
=2\{x=0\}+\{w=0\}$; the exceptional curve is $E=\{x=0\}$ and
$\{w=0\}$ is the strict transform $\widetilde E_0$ of $\{y=0\}$.  Thus the
new special fibre is
\[
 \widetilde E_0+\widetilde E_1+2E ,
\]
and $E$ appears with multiplicity $b_E=2$.  The normalized divisorial
valuation is $v_E=\tfrac12\operatorname{ord}_E$, so
\[
 v_E(x)=\tfrac12,\qquad
 v_E(y)=\tfrac12\operatorname{ord}_E(xw)=\tfrac12 :
\]
the new divisorial point is the midpoint $s=\tfrac12$ of the skeleton.
Blowing up the node subdivides the edge without changing its path metric,
exactly as stated in \cref{sec:running-semistable}.

Now take the vertical divisor $D=E$.  Its local equation is $x$ in chart
(i) and $y$ in chart (ii); on the skeleton this gives
\begin{equation}\label{eq:ch01-node-hat}
 \varphi_E(s)=\min\{s,1-s\},
\end{equation}
the hat function with peak $\tfrac12$ at the midpoint.  The sum of outgoing
slopes gives
\[
 \Delta\varphi_E
 =\delta_{v_0}+\delta_{v_1}-2\delta_{v_{1/2}} .
\]
Intersection theory confirms this with the multiplicity in place.  The
exceptional curve $E$ meets each strict transform transversely in one
point, so $(E\cdot\widetilde E_0)=(E\cdot\widetilde E_1)=1$; and
$\deg\bigl(\cO(\operatorname{div}\pi)|_E\bigr)=0$ gives
$(E\cdot\widetilde E_0)+(E\cdot\widetilde E_1)+2(E\cdot E)=0$, hence
\[
 (E\cdot\widetilde E_0)=(E\cdot\widetilde E_1)=1,
 \qquad
 (E\cdot E)=-1,
\]
and the atomic measure with coefficients
$b_Y\deg(\cO(E)|_Y)$ is
\[
 1\cdot1\cdot\delta_{v_0}
 +1\cdot1\cdot\delta_{v_1}
 +2\cdot(-1)\cdot\delta_{v_{1/2}}
 =\Delta\varphi_E .
\]
Note where the multiplicity $b_E=2$ enters: once in the normalization of
$v_E$, locating the point at the midpoint rather than at $s=1$, and once as
the weight of the atom.  Forgetting either occurrence produces a measure of
the wrong mass; this double bookkeeping is rehearsed again in
\cref{chap:models,chap:metrics}.
\end{example}

This example is the seed of the semistable-reduction picture used
throughout: an equation $xy=a$ contributes an edge of length $v(a)$, models
correspond to subdivisions, and metric-graph potential theory
(\cref{sec:metric-graphs}) computes all curvature.  The nodal degeneration
returns as the basic curve example of the CLD theory in
\cref{chap:cld-examples}, of the BFJ theory in \cref{chap:bfj-examples}, and
of hybrid convergence in \cref{chap:hybrid-spaces}.

\subsection{The toric line}
\label{sec:ch01-running-toric}
\index{toric variety!projective line}

On the dense torus of $\mathbf P^1_K$, the points $v_r$ above form the toric
skeleton $N_{\mathbf R}\simeq\mathbf R$.  A torus-invariant metric on
$\mathcal O(1)$ is encoded, after a fixed reference metric, by a function of
$r$.  Semipositivity becomes convexity with prescribed recession behaviour.
The non-Archimedean Monge--Amp\`ere measure becomes the one-dimensional real
Alexandrov measure, namely the distributional second derivative, with the
normalization established in \cref{chap:model-calculations}.

The precise bookkeeping is as follows; see \cref{sec:toric-convex} for the
general theory and \cite{BGPS14,CLS11} for toric background.  Take the
lattice $N=\mathbf Z$ with dual $M$, and the polytope $P=[0,1]\subset
M_{\mathbf R}$, corresponding to $(\mathbf P^1,\cO(1))$.  Invariant
continuous semipositive metrics on $\cO(1)$ correspond to convex functions
$F:N_{\mathbf R}\to\mathbf R$ with $F-\check h_P$ bounded, where
$\check h_P(u)=\max\{0,-u\}$ is the reflected support function; the
canonical model metric is $F=\check h_P$ itself, with Monge--Amp\`ere
measure $\delta_{\iota(0)}$, the Dirac mass at the Gauss point
(\cref{sec:berkovich-p1}).  Here is a computation with two corners.

\begin{example}[A toric potential with two slope breaks]
\label{ex:ch01-toric-two-breaks}
Let
\[
   F(u)=\max\Bigl\{0,\ -u,\ \tfrac14-\tfrac{u}{2}\Bigr\}.
\]
Comparing the three affine pieces shows that
\[
 F(u)=
 \begin{cases}
  -u,& u\leq-\tfrac12,\\[1mm]
  \tfrac14-\tfrac u2,& -\tfrac12\leq u\leq\tfrac12,\\[1mm]
  0,& u\geq\tfrac12,
 \end{cases}
\]
a convex function with slopes $-1,-\tfrac12,0$ and corners at
$u=\mp\tfrac12$.  Outside $[-\tfrac12,\tfrac12]$ it agrees with
$\check h_P$, and on the middle interval
$0\leq F-\check h_P\leq\tfrac14$, so $F$ encodes a continuous metric on
$\cO(1)$.  Its subgradients at the corners are the intervals
$[-1,-\tfrac12]$ and $[-\tfrac12,0]$, each of length $\tfrac12$, so the real
Monge--Amp\`ere measure of \cref{sec:toric-convex} is
\[
 \operatorname{MA}_{\mathbf R}(F)
 =\tfrac12\delta_{-1/2}+\tfrac12\delta_{1/2},
\]
of total mass $1=\operatorname{vol}_M(P)$.  By
\cref{thm:toric-ma-comparison} with $n=1$ when $F$ is a rational model
metric function, and by \cref{cor:toric-ma-continuous} in general --- the
theorem is stated for a semipositive toric $\mathbf Q$-model metric, which
$F$ need not be here --- the Monge--Amp\`ere measure of the metric is
\[
 c_1\bigl(\cO(1),\|\cdot\|_F\bigr)
 =\tfrac12\delta_{\iota(-1/2)}+\tfrac12\delta_{\iota(1/2)},
\]
two half-atoms at the Gauss points of radii $\rme^{\pm1/2}$.  The function
$F$ always defines a continuous semipositive toric metric.  It is a rational
model metric when
$\tfrac12\in v(K^\times)\otimes_{\mathbf Z}\mathbf Q$; under that
rationality hypothesis a positive multiple is induced by an honest toric
model.  Without it, no model-metric assertion is made.  Rationality issues
of exactly this kind are treated in \cref{chap:metrics}.
\end{example}

The toric line is the one example in which every construction of the book
can be made completely explicit, and it reappears at every stage: for the
CLD calculus in \cref{chap:cld-examples}, for the BFJ theory in
\cref{chap:bfj-examples}, and for the trivially valued theory in
\cref{chap:bj-linear-growth}, where invariant psh functions correspond to
convex functions on $N_{\mathbf R}$ once again.

\subsection{A product test configuration with nontrivial action}
\label{sec:ch01-running-tc}
\index{test configuration!product}

Let $X=\mathbf P^1_k$, $L=\mathcal O(1)$, and let $\mathbf G_m$ act by
\[
                 \lambda\cdot[Z_0:Z_1]=[Z_0:\lambda Z_1].
\]
The product family $X\times\mathbf A^1\to\mathbf A^1$, endowed with this
linearization, is a test configuration although its underlying family is
trivial.  On $H^0(X,mL)$ the monomial $Z_0^{m-i}Z_1^i$ has weight $i$.
Thus geometry of a degeneration can reside entirely in a filtration, even
when the total family is a product.  Its non-Archimedean metric is not
constant: it is a piecewise affine Fubini--Study metric on the valuative cone.

We now compute that metric, its Monge--Amp\`ere measure, and its
Duistermaat--Heckman measure.  The ground field $k$ is algebraically closed
and trivially valued, and we use the notation $v_{a,t}$ of
\cref{sec:ch01-homogeneous-entrance}.

\begin{example}[The Fubini--Study function of the product configuration]
\label{ex:ch01-product-tc}
The filtration of $R_m=H^0(X,mL)$ induced by the action is by weight,
\[
   \mathcal F^\lambda R_m
   =\operatorname{span}\bigl\{Z_0^{m-i}Z_1^i\ \big|\ i\geq\lambda\bigr\},
\]
and the function it defines on $X^{\An}$, in the normalization of
\cref{chap:bj-fs-pl-psh,chap:rays-filtrations}, is
\[
 \phi(v)=\sup_{m\geq1}\ \max_{0\leq i\leq m}
   \frac{i-v\bigl(Z_0^{m-i}Z_1^i\bigr)}{m},
\]
where sections are evaluated through a trivialization of $mL$ at the center
of $v$.  Since $v(Z_0^{m-i}Z_1^i)=(m-i)v(Z_0)+iv(Z_1)$ is affine in $i$,
the inner maximum is attained at $i=0$ or $i=m$, and every $m$ gives the
same value:
\begin{equation}\label{eq:ch01-tc-fs}
 \phi(v)=\max\bigl\{1-v(Z_1),\,-v(Z_0)\bigr\}.
\end{equation}
We evaluate \eqref{eq:ch01-tc-fs} on the cone.  On the ray through
$0=[1:0]$, trivialize with the frame $Z_0$; then $v_{0,t}(Z_0)=0$ and
$v_{0,t}(Z_1)=t$, so
\[
 \phi(v_{0,t})=\max\{1-t,0\}.
\]
On the ray through a point $a\neq0,\infty$ both coordinates are units at
the center, so $\phi(v_{a,t})=\max\{1,0\}=1$.  On the ray through
$\infty=[0:1]$, trivializing with $Z_1$ gives
$\phi(v_{\infty,t})=\max\{1,-t\}=1$.  Thus $\phi$ is the continuous
function equal to $1$ everywhere except on the ray through the fixed point
$0$, where it decreases with slope one until it reaches the value $0$ at
$t=1$ and stays constant beyond.  Subtracting the constant $1$ and
comparing with \cref{ex:ch01-cone-ideal},
\[
 \phi-1=\max\{\varphi_{\mathfrak a},-1\},
\]
the truncation at level $-1$ of the homogeneous potential of the ideal of
the fixed point.

Its Monge--Amp\`ere measure can be computed on the tree by the curve
formula \eqref{eq:bj-curve-laplacian}, proved in \cref{chap:bj-energy}: for
$\deg L=1$,
\[
 \bigl(\omega_{\triv}+\ddc\phi\bigr)
 =\delta_{v_{\triv}}+\Delta\phi ,
\]
where $\Delta$ is the sum-of-outgoing-slopes Laplacian on the tree.  The
function $\phi$ has exactly two corners: at the vertex $v_{\triv}$, where
the single ray through $0$ leaves with outgoing slope $-1$, and at
$v_{0,1}$, where the slopes are $+1$ toward the vertex and $0$ beyond, so
\[
 \Delta\phi=-\delta_{v_{\triv}}+\delta_{v_{0,1}},
 \qquad
 \omega_{\triv}+\ddc\phi=\delta_{v_{0,1}} .
\]
All the curvature sits at the single point $v_{0,1}=\operatorname{ord}_0$:
the non-Archimedean metric of this test configuration is a point mass at a
divisorial valuation, at unit distance along the ray of the attracting
fixed point.
\end{example}

Two asymptotic invariants complete the picture.  The weight of
$\mathbf G_m$ on $H^0(X,mL)$ decomposes with weights $0,1,\ldots,m$, each of
multiplicity one, so the normalized weight measures
\[
   \frac{1}{m+1}\sum_{i=0}^m\delta_{i/m}
\]
converge weakly to Lebesgue measure on $[0,1]$: this is the
Duistermaat--Heckman measure of the configuration, in accordance with
\cref{ex:dict-product-action}; its barycentre $\tfrac12$ is the energy slope
of the associated ray.  Moreover, with $N_m=m+1$ and total weight
$w_m=\tfrac{m(m+1)}2$,
\[
    \frac{w_m}{mN_m}=\frac12
    \qquad\text{for every }m ,
\]
so the $m^{-1}$-term of this expansion vanishes and the Donaldson--Futaki
invariant of the configuration is zero \cite{BHJ17}, as it must be for a
product configuration acting on $\mathbf P^1$.  The configuration is
nevertheless nontrivial: its metric, computed above, is not constant, and
its norm invariants are strictly positive.  Test configurations and their
filtrations are developed in \cref{chap:bj-test-configurations}, their role
in K-stability in \cref{chap:bj-stability} following
\cite{Tia97,Don02,BHJ17}, and the correspondence with geodesic rays and
test curves, which began with \cite{RWN14}, in
\cref{chap:rays-test-curves,chap:rays-filtrations}.

\section{The master comparison diagram}
\label{sec:ch01-master-diagram}
\index{master comparison diagram}

The four entrances produce four kinds of objects.  The later chapters prove
comparison theorems on their stated domains and identify the boundaries
that remain open.  The following diagram is a roadmap whose arrows require
those hypotheses:
{\small
\[
\begin{array}{ccc}
 \text{vertical divisors and model metrics}
   &\overset{\textup{(A)}}{\longrightarrow}& \text{piecewise affine functions on skeleta}\\[2mm]
 \textup{(B)}\big\downarrow && \big\downarrow\textup{(C)}\\[2mm]
 \text{semipositive metrics and global energies}
   &\overset{\textup{(D)}}{\longleftrightarrow}& \text{positive currents and local products}\\[2mm]
 \textup{(E)}\big\uparrow && \big\downarrow\textup{(F)}\\[2mm]
 \text{filtrations, rays, and test configurations}
   &\underset{\textup{(G)}}{\longrightarrow}& \text{Monge--Amp\`ere and slope measures}.
\end{array}
\]
}
The upper horizontal arrow is an identification on a fixed suitable model.
The two middle theories agree on semipositive model metrics and, under stated
regularizability hypotheses, on certain continuous limits.  The lower left
arrow is realized by Rees and Fubini--Study constructions in the upward
direction and by asymptotic slopes of complex energies in the downward one.
No unrestricted equivalence of all local CLD-psh functions
with all global BFJ/BJ-psh metrics is asserted.

Each arrow has a home in the book, and it is worth recording them now.

\begin{enumerate}
\item[(A)] A vertical divisor on a model restricts to a piecewise affine
function on the skeleton of that model, integral-affine on each face, and
every rational piecewise affine function on a skeleton arises this way
after subdivision.  This is proved in \cref{chap:skeleta}
(\cref{prop:skeleton-model-pa,thm:skeleton-pl-realization}), and
\cref{ex:ch01-blowup-model-function,ex:ch01-node-blowup} are its two
smallest cases.

\item[(B)] Passing from model metrics to their uniform limits and then to
decreasing singular limits produces the continuous semipositive and
singular psh metrics of the global theory: over a nontrivially valued field
this is the BFJ theory of \cref{part:bfj}
(\cref{chap:bfj-psh,chap:bfj-ma}), and over a trivially valued field the BJ
theory of \cref{part:bj}.  The arrow preserves Monge--Amp\`ere mass along
its uniform-limit half, by \cref{thm:cl-continuous-extension}, and for
bounded psh potentials by \cref{prop:bounded-ma-no-polar-mass}; on
$\cE^1$ it does so with the nonpluripolar convention
(\cref{def:nonpluripolar-ma}).  For a general singular psh metric mass can be
lost to the polar set, and the arrow is then only mass-nonincreasing.

\item[(C)] A piecewise affine function on a skeleton, viewed through
tropical charts, has a curvature current in the Chambert--Loir--Ducros
calculus, and \emph{continuous} psh functions admit Bedford--Taylor products
(not merely locally bounded ones: see \cref{chap:cld-bedford-taylor}).
This is the content of \cref{part:cld}, especially
\cref{chap:cld-bedford-taylor,chap:cld-metrized-bundles};
\cref{ex:ch01-lagerberg-corner} is the one-dimensional germ.

\item[(D)] The two middle theories are compared in
\cref{chap:compare-ma}: on semipositive model metrics the identification
is \cref{prop:model-core-preview} below; on continuous semipositive metrics
it passes through uniform approximation; for singular metrics only partial
statements are true, and \cref{chap:master-dictionary} records the exact
status of each locus.

\item[(E)] Test configurations, filtrations, and geodesic rays produce
semipositive metrics and, under the stated energy hypotheses, finite-energy
objects through the Rees and Fubini--Study constructions of
\cref{chap:bj-test-configurations,chap:rays-filtrations};
\cref{ex:ch01-product-tc} is the smallest nontrivial instance.  On the
complex side the same arrow is the Ross--Witt~Nystr\"om correspondence of
\cref{chap:rays-test-curves}.

\item[(F)] Positive currents and metrics have Monge--Amp\`ere measures;
their comparison across the four theories, including total masses and the
model core, is the subject of \cref{chap:compare-ma}.

\item[(G)] Filtrations have Duistermaat--Heckman and slope measures, which
match Monge--Amp\`ere data through energy asymptotics; these comparisons
are proved in \cref{chap:rays-filtrations,chap:transcendental-okounkov-dh},
and the vanishing computed in \cref{sec:ch01-running-tc} is the simplest
slope identity.
\end{enumerate}

The complex and transcendental theories of \cref{part:complex} attach to a
compact K\"ahler \emph{manifold} a space of non-Archimedean metrics
(\cref{chap:transcendental-metrics,chap:mp-psh}; the extension to singular
K\"ahler spaces is cited, not proved --- see \cref{rem:tna-smooth-scope});
their comparison with the
trivially valued column of the diagram is the order-preserving bijection of
\cref{chap:dxz-bj-comparison}, and the passage from a complex degeneration
to its non-Archimedean limit is realized by the hybrid spaces of
\cref{chap:hybrid-spaces}.  The full diagram, with every arrow typed as
equivalence, injection, one-way functor, or analogy, is stated and proved
as \cref{thm:dict-master-diagram} in \cref{chap:master-dictionary}.  All
three running examples traverse the entire diagram, and the reader is
encouraged to trace them through it: the computations of
\cref{sec:ch01-running-examples} already populate five of its seven arrows.

On the common model core, the comparison is not merely an analogy, and we
can state it precisely now.

\begin{proposition}\label{prop:model-core-preview}
Assume that $K$ is discretely valued.  Let $X$ be a proper, normal,
pure $n$-dimensional $K$-scheme and let $L_1,\ldots,L_n$ carry semipositive
model metrics on one common proper flat normal model, represented by line
bundles extending $L_i^{\otimes m_i}$ for positive integers $m_i$.  For the
associated rational formal metrics, the CLD mixed top-degree product agrees
with the Chambert--Loir intersection measure.
\end{proposition}
\begin{proof}
We spell out the reduction because it fixes the normalization used in both
parts of the book.  Replace the common algebraic model by its formal
completion and then by a common admissible formal blowup on which all the
metrics are represented by line bundles
$\mathfrak L_1,\ldots,\mathfrak L_n$, with the same exponents $m_i$.
The first Chern current of the induced rational formal metric on $L_i$ is
$m_i^{-1}c_1(\mathfrak L_i)$.  Pullback changes neither side: on the
intersection-theoretic side this is the projection formula, and on the CLD
side it is compatibility of first Chern currents with proper pullback and
pushforward.

Write the special fibre as $\sum_Y b_Y Y$, where $Y$ runs over its
irreducible components, and let $x_Y$ be the corresponding divisorial point.
The Chambert--Loir measure is, by definition,
\begin{equation}\label{eq:preview-cl-measure}
 \frac{1}{m_1\cdots m_n}\sum_Y b_Y
 \deg_Y\!\left(
   c_1(\mathfrak L_1|_Y)\cdots c_1(\mathfrak L_n|_Y)\cap[Y]
 \right)\delta_{x_Y}.
\end{equation}
The residue-field degree is included in the displayed degree.

We now compute the CLD product.  Choose, after replacing the
$\mathfrak L_i$ by sufficiently positive differences and using
multilinearity, meromorphic sections whose divisors meet the strata of the
special fibre properly.  The Poincar\'e--Lelong formula replaces the first
Chern current of each formal metric by the divisor current of its section
plus an exact $\dprimeddprime$-term.  Apply this successively.  Stokes' theorem
removes the exact terms, while intersection with the divisor current
restricts the remaining formal line bundles to that divisor.  Induction on
the dimension therefore reduces the coefficient over the generic point of
$Y$ to
\[
 \frac{b_Y}{m_1\cdots m_n}\deg_Y\!\left(
   c_1(\mathfrak L_1|_Y)\cdots c_1(\mathfrak L_n|_Y)\cap[Y]
 \right).
\]
There is no contribution away from the divisorial points: on a formal chart
the model potentials are integer-affine on every open face, so their top
corner product is supported where the reduction has a top-dimensional
component.  The CLD product is consequently exactly the atomic measure in
\eqref{eq:preview-cl-measure}.

The two foundational ingredients just used---the formal-metric
Poincar\'e--Lelong identity and the identification of the residual local
degree with the calibrated tropical multiplicity---are proved in
\cite[Theorem~18.5.5 and Proposition~18.5.8]{CLD25}.  They are reconstructed
in \Cref{chap:cld-metrized-bundles}; the argument above is the induction in
that source, with all normalization factors displayed.
\end{proof}

The proposition asserts an equality of measures on the model core only.  It
does not assert that a singular BFJ or BJ psh metric has a CLD wedge
product, nor that every CLD-psh function is a limit of model functions;
those finer questions have partial answers whose exact boundaries are the
business of \cref{chap:compare-ma} and of the status ledger in
\cref{chap:master-dictionary}.

\section{How to read the book}
\label{sec:ch01-readers-guide}
\index{reading routes}

\Crefrange{chap:valuations}{chap:model-calculations} constructs the common
language.  \Cref{part:cld} develops the local differential theory in a
concrete-to-invariant order.  \Cref{part:bfj} develops the BFJ global
variational theory over nontrivially valued fields.  \Cref{part:bj} treats
the trivially valued theory, including the full envelope theorem for
unibranch projective varieties.  \Cref{part:complex} constructs the complex
and transcendental non-Archimedean spaces.  \Cref{part:comparisons} proves
the comparison arrows and records their exact domains.

\subsection{What each part proves}
\label{sec:ch01-part-summaries}

\begin{itemize}
\item \Cref{part:foundations} constructs valuations, models, Berkovich
spaces, skeleta, and metrics, and ends with the computational test suite of
\cref{chap:model-calculations}.  Its headline results are the
skeleton/retraction theorems of \cref{chap:skeleta} and the model-measure
theory of \cref{chap:metrics}.  Everything later depends on it.

\item \Cref{part:cld} builds the Chambert--Loir--Ducros calculus from
Lagerberg superforms through tropicalization to currents, the
Poincar\'e--Lelong formula, and Bedford--Taylor products
(\cref{chap:cld-bedford-taylor}), together with metrized bundles and Chern
currents (\cref{chap:cld-metrized-bundles}) and the delta-form and
harmonic-tropicalization extensions (\cref{chap:delta-forms}).

\item \Cref{part:bfj} develops the global theory over a nontrivially valued
field: compactness and regularization of psh functions, envelopes, the
Monge--Amp\`ere operator, capacity, energy, and the variational solution of
the non-Archimedean Monge--Amp\`ere equation, with the original
hypotheses of \cite{BFJ16,BFJ15} and their later extensions kept visibly
separate (\cref{chap:bfj-continuity}).

\item \Cref{part:bj} develops the trivially valued theory: the cone of
semivaluations, test configurations and filtrations, Fubini--Study and psh
functions, homogeneity and b-divisors, and then the complete proof of the
envelope theorem for integral projective unibranch varieties in arbitrary
characteristic (\cref{chap:bj-envelope-statement}, proved in
\crefrange{chap:bj-cartier-realization}{chap:bj-envelope-consequences}),
followed by energy, finite-energy measures, and stability functionals.

\item \Cref{part:complex} constructs non-Archimedean geometry from complex
geometry: test curves and geodesic rays, the transcendental metrics of
Darvas--Xia--Zhang (\cref{chap:transcendental-metrics}), partial Okounkov
bodies and Duistermaat--Heckman measures, the comparison with
Boucksom--Jonsson (\cref{chap:dxz-bj-comparison}), the tropical spectra and
psh theory of Mesquita-Piccione (\cref{chap:mp-psh}), and the
transcendental Calabi--Yau theorem with its cscK applications
(\cref{chap:mpwn-cy,chap:mpwn-k-stability}).

\item \Cref{part:comparisons} proves the master dictionary
(\cref{chap:master-dictionary}), compares the Monge--Amp\`ere operators
(\cref{chap:compare-ma}), identifies rays, filtrations, and Legendre
transforms (\cref{chap:rays-filtrations}), develops hybrid convergence
(\cref{chap:hybrid-spaces}), functoriality and trace
(\cref{chap:functoriality-trace}), and the Yau--Tian--Donaldson bridge
(\cref{chap:ytd-bridge}), and presents verified extensions and structural
obstructions in \cref{chap:verified-extensions}.  Its final three chapters
develop local valuative psh functions
(\cref{chap:local-valuative-psh}), the pseudoeffective DXZ--MP comparison
(\cref{chap:pseff-mp-dxz}), and the remaining research frontiers
(\cref{chap:major-open-problems}).
\end{itemize}

\subsection{Three reading routes}
\label{sec:ch01-reading-routes}

The book is long, and three coherent routes traverse it.

\paragraph{The local differential route.}
Read \cref{part:foundations} and then all of \cref{part:cld}.  Continue
with the continuity theory of \cref{chap:bfj-continuity}, the envelope
consequences of \cref{chap:bj-envelope-consequences}, the transcendental
Calabi--Yau theorem of \cref{chap:mpwn-cy}, and conclude with the
Monge--Amp\`ere comparisons of \cref{chap:compare-ma}.  This route suits a
reader whose interests are local: forms, currents, and the analysis of the
operators themselves.

\paragraph{The global variational route.}
Read \cref{part:foundations}, then \cref{part:bfj,part:bj} in order, and
finish with \crefrange{chap:mp-psh}{chap:hybrid-spaces}.  On this route the
CLD calculus of \cref{part:cld} may be treated as a black box whose outputs
are quoted; the emphasis is on compactness, envelopes, energy, and the
variational method, culminating in the envelope theorem and the
non-Archimedean Calabi--Yau theorems.

\paragraph{The complex/K\"ahler route.}
Read \crefrange{chap:four-entrances}{chap:metrics}, consult
\cref{chap:topological-tools} for the topological and convex tools as needed, and work
through the examples of \cref{chap:model-calculations}.  Then read
\cref{part:complex} and
\crefrange{chap:master-dictionary}{chap:hybrid-spaces}.
This route is designed for a complex geometer who wants the
non-Archimedean limit theory of degenerating K\"ahler metrics and the
cscK/K-stability applications; the deeper complex pluripotential
prerequisites are restated at their interfaces in
\cref{chap:complex-prerequisites}, with detailed background in
\cite{XiaBook}.

\medskip

Whatever the route, \cref{chap:model-calculations} should not be skipped:
it is short, entirely computational, and every later chapter calibrates its
signs against it.

\subsection{A dictionary of first glances}
\label{sec:ch01-first-glance-dictionary}
\index{dictionary!first-glance}

The following table lists, informally, the correspondences that the book
makes precise.  Each row is an analogy at this point; the right column
records where it becomes a construction or a theorem.

\begin{center}
\small
\begin{tabular}{p{.31\textwidth}|p{.33\textwidth}|p{.24\textwidth}}
\textbf{Complex geometry}&\textbf{Non-Archimedean geometry}
  &\textbf{Where made precise}\\ \hline
smooth Hermitian metric
  &model or PL metric from a model
  &\cref{chap:models,chap:metrics}\\
psh potential
  &BFJ/BJ psh function; CLD psh function
  &\cref{chap:bfj-psh,chap:bj-fs-pl-psh,chap:cld-poincare-lelong}\\
$\ddc$
  &$\dprimeddprime$ on tropical charts; slope Laplacian on skeleta
  &\cref{chap:superforms,chap:model-calculations}\\
Bedford--Taylor product
  &Chambert--Loir measure; CLD wedge; BFJ/BJ Monge--Amp\`ere
  &\cref{chap:cld-bedford-taylor,chap:bfj-ma,chap:bj-energy}\\
equilibrium measure of the unit disc
  &Dirac mass at the Gauss point
  &\cref{sec:ch01-equilibrium,sec:berkovich-p1}\\
Lelong numbers, multiplier ideals
  &values on divisorial valuations
  &\cref{chap:valuative-analytification}\\
geodesic ray, test curve
  &test configuration, filtration
  &\cref{chap:rays-test-curves,chap:rays-filtrations}\\
Mabuchi/Ding functionals
  &non-Archimedean stability functionals
  &\cref{chap:bj-stability,chap:ytd-bridge}\\
degenerating family over the disc
  &hybrid space; non-Archimedean central fibre
  &\cref{chap:hybrid-spaces}
\end{tabular}
\end{center}

A warning accompanies the table.  Rows are not bijections: the two sides of
each row have different regularity classes, different compactness
properties, and sometimes different spaces underneath.  The precise
comparison statements, with hypotheses, are the content of
\cref{part:comparisons}, and no row of this table is ever cited as a
substitute for them.

\section*{Source notes}
\label{sec:ch01-source-notes}

The analytic and model viewpoints originate in Berkovich's analytic geometry
\cite{Berk12}, Zhang's semipositive metrics \cite{Zhang1995}, and the
work of Chambert--Loir \cite{CL06}. The global singular theory is
developed in \cite{BFJ16,BFJ15,BFJ16b}; the trivially valued theory in
\cite{BJ22}; and the differential theory in
\cite{CLD12}, whose revised version \cite{CLD25} is our reference text for
\cref{part:cld}.  Superforms were introduced by Lagerberg in 2012; his
paper, together with Thuillier's earlier potential theory on Berkovich
curves and the Baker--Rumely theory of the Berkovich projective line, is
part of the prehistory of the differential entrance and is cited at the
relevant points of \cref{part:cld} and \cref{chap:model-calculations}.

The complex model-case computation of \cref{sec:ch01-equilibrium} is
classical logarithmic potential theory; the global framework into which it
fits is Bedford--Taylor theory \cite{BT76,BT82} and its finite-energy and
variational extensions \cite{BEGZ10,BBGZ13,GZ17}.  The K-stability context
for test configurations begins with \cite{Tia97,Don02}; the
Duistermaat--Heckman and norm formalism used for the running example is
from \cite{BHJ17}; and the complex-to-non-Archimedean constructions
compared in \cref{part:complex,part:comparisons} are those of
\cite{RWN14,DXZ23,MP24,MPWN25}, with the lecture notes \cite{BJ26ComplexNA}
and the hybrid-space methods of \cite{BJ17Hybrid} providing the bridge
narrative.  Toric background is \cite{CLS11,BGPS14}.

This chapter's expository order is our own.
In particular, \cref{part:cld} does not follow the order of the CLD manuscript.
